\documentclass[exercice,a4]{thedocument}%
\usepackage{pstricks}%
\usepackage{tabularx}%
\usepackage[np]{numprint}%
\startdatakeys
    \source{25GENMATNC1}
    \BO{2026}
    \cycle{4}
    \competencesDuSocle{CA,RA}
    \objectifsApprentissage{C4A1aa,C4A2ad,C4B2cb,C4D2aa,C4D3ac}
    \motsCles{nombre premier,aire,fonction affine,vitesse,pourcentage}
\enddatakeys
\begin{document}%
\providecommand{\degres}{^\circ}%
    

\medskip

Cet exercice est un questionnaire à choix multiples (QCM).

Pour chaque question, une seule des trois réponses proposées est exacte.

Sur la copie, indiquer le numéro de la question et la réponse A, B ou C choisie. Aucune justification n'est demandée.

Aucun point ne sera enlevé en cas de mauvaise réponse.

\begin{center}
\begin{tabularx}{\linewidth}{|c|m{5.5cm}|*{3}{>{\centering \arraybackslash}X|}}\hline
\multicolumn{2}{|c|}{Questions}						&Réponse A	&Réponse B	&Réponse C\\ \hline
\textbf{1.}&Parmi les nombres suivants, lequel est premier ?	&719		&935		&687\\ \hline
\textbf{2.}&Quelle est l'aire de la figure ci-dessous ?

\begin{pspicture}(5.6,5)
\psset{arrowsize=2pt 3}
\def\bb{\psline(-0.05,-0.1)(-0.05,0.1)\psline(0.05,-0.1)(0.05,0.1)}
\pspolygon(0,4)(4.6,4)(4.6,0.7)(3.3,0.7)(3.3,2)(0,2)
\psframe(0,4)(0.2,3.8)\psframe(0,2)(0.2,2.2)\psframe(4.6,4)(4.4,3.8)
\psframe(4.6,0.7)(4.4,0.9)\psframe(3.3,0.7)(3.5,0.9)
\rput(3.95,0.7){\bb}\rput{90}(3.3,1.35){\bb}
\psline[linewidth=0.5pt]{<->}(0,4.2)(4.6,4.2)\uput[u](2.3,4.2){\small 7 cm}
\psline[linewidth=0.5pt]{<->}(4.8,4)(4.8,0.7)\uput[r](4.8,2.2){\small 5 cm}
\psline[linewidth=0.5pt]{<->}(4.6,0.5)(3.3,0.5)\uput[d](3.95,0.5){\small 2 cm}
\end{pspicture}&24 cm$^2$&140 cm$^2$&25 cm$^2$\\ 
\hline
\textbf{3.}& Une de ces fonctions est affine.\newline Laquelle ?&$f$ définie par $f(x) = 3 (x + 1 )$&$g$ définie par $g(x) = \dfrac{5}{x} + 1$\rule[-10pt]{0mm}{0pt}&$h$ définie par $h(x) =  x^2 + 1$\\ 
\hline
\textbf{4.}&La distance de Tontouta à Narita est égale à environ 6 980 km. Le vol Tontouta-Narita dure environ 9 heures. Quelle est la vitesse moyenne, arrondie à la centaine de km/h, de 
l'avion sur ce trajet ?&600 km/h&800 km/h&\np{1000} km/h\\
 \hline
\textbf{5.}&Dans un collège de 730 élèves, 60\,\% des élèves sont des filles.

Quel est le nombre de filles dans ce collège ?&438&60&670\\ 
\hline
\end{tabularx}
\end{center}

\newpage
\showthesource
\end{document}
